%!TEX root = ../main.tex
% Chapter 8: Generative Models --- GAN
\chapter{Generative Models --- GAN}
\label{ch:gan}

\paragraph{Chapter objectives}
\begin{itemize}
\item Distinguish generative and discriminative models
\item Derive the minimax objective of GANs and the optimal discriminator
\item Understand training problems and solutions (WGAN, WGAN-GP)
\item Implement a complete DCGAN in PyTorch
\item Know the evaluation metrics (FID, IS)
\end{itemize}

\section{Generative vs discriminative models}
\label{sec:gen-vs-disc}

\begin{definition}[Discriminative model\index{discriminative model}]
A discriminative model directly learns the conditional distribution $p(y \mid x)$, i.e., the decision boundary between classes. Examples: logistic regression, SVM, classifier neural networks.
\end{definition}

\begin{definition}[Generative model\index{generative model}]
A generative model learns the joint distribution $p(x, y)$ or the marginal distribution $p(x)$ of the data. It can \emph{generate} new samples. Examples: GAN, VAE (cf.\ Chapter~\ref{ch:vae}), diffusion models, autoregressive models.
\end{definition}

\begin{remark}
A generative model not only allows creating new data, but also performing probabilistic reasoning (anomaly detection, completion, etc.).
\end{remark}

\section{GAN framework: the minimax objective}
\label{sec:gan-framework}

\subsection{Formulation}

A \textbf{Generative Adversarial Network}\index{GAN} (Goodfellow et al., 2014) consists of two networks in competition:
\begin{itemize}
    \item The \textbf{generator} $G : \mathcal{Z} \to \mathcal{X}$ transforms random noise $z \sim p_z(z)$ into a sample $G(z)$.
    \item The \textbf{discriminator} $D : \mathcal{X} \to [0, 1]$ estimates the probability that a sample comes from the real data rather than from the generator.
\end{itemize}

\begin{mathbox}[title=\textbf{GAN minimax objective}]\index{minimax objective}
\begin{equation}
\min_G \max_D \; V(D, G) = \E_{x \sim p_{\text{data}}}\!\big[\log D(x)\big] + \E_{z \sim p_z}\!\big[\log\!\big(1 - D(G(z))\big)\big]
\label{eq:gan-minimax}
\end{equation}
\end{mathbox}

\subsection{Derivation of the optimal discriminator $D^*$}

\begin{theorem}[Optimal discriminator\index{optimal discriminator}]
\label{thm:optimal-D}
For a fixed generator $G$, the optimal discriminator is:
\begin{equation}
D^*_G(x) = \frac{p_{\text{data}}(x)}{p_{\text{data}}(x) + p_g(x)}
\label{eq:optimal-D}
\end{equation}
where $p_g$ is the distribution induced by $G$.
\end{theorem}

\begin{proof}
For fixed $G$, we maximize $V(D, G)$ with respect to $D$. The integrand is, for each $x$:
\begin{equation}
f(D(x)) = p_{\text{data}}(x) \log D(x) + p_g(x) \log(1 - D(x))
\end{equation}
This is a function of $D(x) \in [0, 1]$. Differentiating and setting to zero:
\begin{align}
\frac{\partial f}{\partial D(x)} &= \frac{p_{\text{data}}(x)}{D(x)} - \frac{p_g(x)}{1 - D(x)} = 0 \\
\Rightarrow \quad p_{\text{data}}(x)\big(1 - D(x)\big) &= p_g(x) D(x) \\
\Rightarrow \quad D^*(x) &= \frac{p_{\text{data}}(x)}{p_{\text{data}}(x) + p_g(x)}
\end{align}
One verifies that the second derivative is negative, confirming the maximum.
\end{proof}

\subsection{Derivation of generator optimality}

\begin{theorem}[Generator optimality and Jensen-Shannon divergence\index{Jensen-Shannon divergence}]
\label{thm:optimal-G}
Substituting $D^*$ into the objective, we obtain:
\begin{equation}
C(G) = V(D^*_G, G) = -\log 4 + 2 \cdot \KL{JSD}(p_{\text{data}} \| p_g)
\end{equation}
where the Jensen-Shannon divergence is:
\begin{equation}
\text{JSD}(p \| q) = \frac{1}{2}\KL(p \| m) + \frac{1}{2}\KL(q \| m), \quad m = \frac{p + q}{2}
\end{equation}
The global minimum $C(G) = -\log 4$ is achieved if and only if $p_g = p_{\text{data}}$.
\end{theorem}

\begin{proof}
Substituting $D^*$:
\begin{align}
C(G) &= \E_{x \sim p_{\text{data}}}\!\left[\log \frac{p_{\text{data}}(x)}{p_{\text{data}}(x) + p_g(x)}\right] + \E_{x \sim p_g}\!\left[\log \frac{p_g(x)}{p_{\text{data}}(x) + p_g(x)}\right] \\
&= \E_{p_{\text{data}}}\!\left[\log \frac{p_{\text{data}}}{2 \cdot \frac{p_{\text{data}} + p_g}{2}}\right] + \E_{p_g}\!\left[\log \frac{p_g}{2 \cdot \frac{p_{\text{data}} + p_g}{2}}\right] \\
&= -\log 4 + \KL\!\left(p_{\text{data}} \,\Big\|\, \frac{p_{\text{data}} + p_g}{2}\right) + \KL\!\left(p_g \,\Big\|\, \frac{p_{\text{data}} + p_g}{2}\right) \\
&= -\log 4 + 2 \cdot \text{JSD}(p_{\text{data}} \| p_g)
\end{align}
Since $\text{JSD} \geq 0$ with equality iff $p_{\text{data}} = p_g$, the minimum is indeed $-\log 4$.
\end{proof}

\subsection{TikZ diagram of GAN training}

\begin{figure}[htbp]
\centering
\begin{tikzpicture}[
    netblock/.style={rectangle, draw, fill=blue!12, minimum width=2.8cm, minimum height=1.2cm, align=center, rounded corners=3pt, font=\small\bfseries},
    datablock/.style={rectangle, draw, fill=green!12, minimum width=2.2cm, minimum height=0.9cm, align=center, rounded corners=3pt, font=\small},
    lossblock/.style={rectangle, draw, fill=red!12, minimum width=2.5cm, minimum height=0.9cm, align=center, rounded corners=3pt, font=\small},
    arr/.style={-{Stealth}, thick},
    darr/.style={-{Stealth}, thick, dashed},
]

% Generator
\node[datablock] (noise) at (0, 0) {$z \sim \mathcal{N}(0, I)$};
\node[netblock] (G) at (3, 0) {Generator\\$G_\theta$};
\node[datablock] (fake) at (6.5, 0) {$G(z)$\\(fake)};

% Real data
\node[datablock] (real) at (6.5, 2.5) {$x \sim p_{\text{data}}$\\(real)};

% Discriminator
\node[netblock] (D) at (10, 1.25) {Discriminator\\$D_\phi$};

% Output
\node[lossblock] (loss) at (13.5, 1.25) {$V(D, G)$\\Loss};

% Arrows
\draw[arr] (noise) -- (G);
\draw[arr] (G) -- (fake);
\draw[arr] (fake) -- (D);
\draw[arr] (real) -- (D);
\draw[arr] (D) -- (loss);

% Backpropagation
\draw[darr, red!60!black] (loss.south) -- ++(0, -1.2) -| (D.south)
    node[pos=0.25, below, font=\footnotesize, text=red!60!black] {$\nabla_\phi$ (maximize $V$)};
\draw[darr, blue!60!black] (loss.north) -- ++(0, 1.0) -| (G.north)
    node[pos=0.25, above, font=\footnotesize, text=blue!60!black] {$\nabla_\theta$ (minimize $V$)};

\end{tikzpicture}
\caption{GAN training loop. The discriminator $D$ is updated to maximize $V$, while the generator $G$ is updated to minimize it. Dashed arrows represent backpropagation.}
\label{fig:gan-training}
\end{figure}

\section{Training problems}
\label{sec:gan-training-problems}

\subsection{Mode collapse}

\begin{definition}[Mode collapse\index{mode collapse}]
The generator ``specializes'' in producing a small number of modes of the target distribution, ignoring the diversity of the real data. Formally, $p_g$ is concentrated on a subset of small measure within the support of $p_{\text{data}}$.
\end{definition}

\subsection{Vanishing gradients for $G$}

\begin{warning}
If the discriminator is too powerful ($D(x) \approx 1$ for real data and $D(G(z)) \approx 0$ for fake data), then $\log(1 - D(G(z))) \approx \log(1) = 0$ and the gradient for $G$ becomes near zero.

\textbf{Practical solution:} instead of minimizing $\log(1 - D(G(z)))$, maximize $\log D(G(z))$ (``non-saturating loss''\index{non-saturating loss}).
\end{warning}

\section{Wasserstein GAN (WGAN)}
\label{sec:wgan}

\subsection{Earth Mover's distance (Wasserstein-1)}

\begin{definition}[Wasserstein-1 distance\index{Wasserstein distance}\index{Earth Mover's distance}]
\begin{equation}
W(p_{\text{data}}, p_g) = \inf_{\gamma \in \Pi(p_{\text{data}}, p_g)} \E_{(x, y) \sim \gamma}\!\big[\|x - y\|\big]
\end{equation}
where $\Pi(p_{\text{data}}, p_g)$ is the set of joint distributions whose marginals are $p_{\text{data}}$ and $p_g$.
\end{definition}

\begin{theorem}[Kantorovich-Rubinstein duality]
\begin{equation}
W(p_{\text{data}}, p_g) = \sup_{\|f\|_L \leq 1} \E_{x \sim p_{\text{data}}}\!\big[f(x)\big] - \E_{x \sim p_g}\!\big[f(x)\big]
\label{eq:kantorovich}
\end{equation}
where the supremum is taken over 1-Lipschitz\index{Lipschitz} functions.
\end{theorem}

\begin{intuition}
The Wasserstein distance is smoother than the JSD: even when the supports of $p_{\text{data}}$ and $p_g$ do not overlap, $W$ provides a usable gradient, whereas $\KL$ and JSD are either infinite or constant.
\end{intuition}

\begin{keyformulas}[title=\textbf{WGAN objective}]\index{WGAN}
The discriminator becomes a ``critic'' $f_w$ (without a sigmoid at the output):
\begin{equation}
\max_{w : \|f_w\|_L \leq 1} \; \E_{x \sim p_{\text{data}}}\!\big[f_w(x)\big] - \E_{z \sim p_z}\!\big[f_w(G_\theta(z))\big]
\label{eq:wgan-objective}
\end{equation}
\end{keyformulas}

\subsection{WGAN-GP: gradient penalty}
\label{sec:wgan-gp}

Arjovsky et al.\ enforce the Lipschitz constraint through \emph{weight clipping}, which is problematic. Gulrajani et al.\ (2017) propose a \textbf{gradient penalty}\index{gradient penalty}\index{WGAN-GP}:

\begin{keyformulas}[title=\textbf{WGAN-GP}]
\begin{equation}
\mathcal{L}_{\text{WGAN-GP}} = \underbrace{\E_{z}\!\big[f_w(G(z))\big] - \E_{x}\!\big[f_w(x)\big]}_{\text{WGAN objective}} + \lambda \,\underbrace{\E_{\hat{x}}\!\big[\big(\|\nabla_{\hat{x}} f_w(\hat{x})\|_2 - 1\big)^2\big]}_{\text{gradient penalty}}
\label{eq:wgan-gp}
\end{equation}
where $\hat{x} = \epsilon x + (1 - \epsilon) G(z)$ with $\epsilon \sim U(0, 1)$ and typically $\lambda = 10$.
\end{keyformulas}

\section{Conditional GAN (cGAN)}
\label{sec:cgan}

\begin{definition}[Conditional GAN\index{conditional GAN}\index{cGAN}]
Both $G$ and $D$ are conditioned on auxiliary information $y$ (class label, text, etc.):
\begin{equation}
\min_G \max_D \; \E_{x, y}\!\big[\log D(x, y)\big] + \E_{z, y}\!\big[\log\big(1 - D(G(z, y), y)\big)\big]
\end{equation}
\end{definition}

\section{DCGAN: architectural guidelines}
\label{sec:dcgan}

The \textbf{Deep Convolutional GAN}\index{DCGAN} (Radford et al., 2016) establishes empirical rules to stabilize training of convolutional GANs:

\begin{bestpractice}
\begin{enumerate}
    \item Replace pooling with \emph{strided convolutions} in $D$ and transposed convolutions in $G$.
    \item Use \emph{Batch Normalization}\index{Batch Normalization} in $G$ and $D$ (except the last layer of $G$ and the first layer of $D$).
    \item Remove fully connected layers.
    \item Use ReLU in $G$ (except the output: $\tanh$) and LeakyReLU in $D$.
\end{enumerate}
\end{bestpractice}

\section{PyTorch implementation: DCGAN on Fashion-MNIST}
\label{sec:dcgan-pytorch}

\begin{codeblock}[title=\textbf{Complete DCGAN with training loop}]
\begin{minted}{python}
import torch
import torch.nn as nn
from torchvision import datasets, transforms
from torch.utils.data import DataLoader

# Hyperparameters
LATENT_DIM = 100
IMG_CHANNELS = 1
FEATURES_G = 64
FEATURES_D = 64
LR = 2e-4
BETAS = (0.5, 0.999)
BATCH_SIZE = 128
NUM_EPOCHS = 25

class Generator(nn.Module):
    """DCGAN generator for 28x28 images."""
    def __init__(self, latent_dim, features_g, img_channels):
        super().__init__()
        self.net = nn.Sequential(
            # Input: (batch, latent_dim, 1, 1)
            nn.ConvTranspose2d(latent_dim, features_g * 4, 4, 1, 0,
                               bias=False),
            nn.BatchNorm2d(features_g * 4),
            nn.ReLU(True),
            # -> (batch, features_g*4, 4, 4)
            nn.ConvTranspose2d(features_g * 4, features_g * 2, 3, 2, 1,
                               bias=False),
            nn.BatchNorm2d(features_g * 2),
            nn.ReLU(True),
            # -> (batch, features_g*2, 7, 7)
            nn.ConvTranspose2d(features_g * 2, features_g, 4, 2, 1,
                               bias=False),
            nn.BatchNorm2d(features_g),
            nn.ReLU(True),
            # -> (batch, features_g, 14, 14)
            nn.ConvTranspose2d(features_g, img_channels, 4, 2, 1,
                               bias=False),
            nn.Tanh(),
            # -> (batch, img_channels, 28, 28)
        )

    def forward(self, z):
        return self.net(z)


class Discriminator(nn.Module):
    """DCGAN discriminator for 28x28 images."""
    def __init__(self, img_channels, features_d):
        super().__init__()
        self.net = nn.Sequential(
            # Input: (batch, img_channels, 28, 28)
            nn.Conv2d(img_channels, features_d, 4, 2, 1, bias=False),
            nn.LeakyReLU(0.2, inplace=True),
            # -> (batch, features_d, 14, 14)
            nn.Conv2d(features_d, features_d * 2, 4, 2, 1, bias=False),
            nn.BatchNorm2d(features_d * 2),
            nn.LeakyReLU(0.2, inplace=True),
            # -> (batch, features_d*2, 7, 7)
            nn.Conv2d(features_d * 2, features_d * 4, 3, 2, 1, bias=False),
            nn.BatchNorm2d(features_d * 4),
            nn.LeakyReLU(0.2, inplace=True),
            # -> (batch, features_d*4, 4, 4)
            nn.Conv2d(features_d * 4, 1, 4, 1, 0, bias=False),
            nn.Sigmoid(),
            # -> (batch, 1, 1, 1)
        )

    def forward(self, x):
        return self.net(x).view(-1, 1)


def weights_init(m):
    """Weight initialization following DCGAN recommendations."""
    classname = m.__class__.__name__
    if 'Conv' in classname:
        nn.init.normal_(m.weight.data, 0.0, 0.02)
    elif 'BatchNorm' in classname:
        nn.init.normal_(m.weight.data, 1.0, 0.02)
        nn.init.constant_(m.bias.data, 0)


# --- Data preparation ---
transform = transforms.Compose([
    transforms.ToTensor(),
    transforms.Normalize([0.5], [0.5]),  # -> [-1, 1]
])
dataset = datasets.FashionMNIST(
    root='./data', train=True, download=True, transform=transform
)
dataloader = DataLoader(dataset, batch_size=BATCH_SIZE, shuffle=True,
                        num_workers=2, drop_last=True)

# --- Initialization ---
device = torch.device('cuda' if torch.cuda.is_available() else 'cpu')
G = Generator(LATENT_DIM, FEATURES_G, IMG_CHANNELS).to(device)
D = Discriminator(IMG_CHANNELS, FEATURES_D).to(device)
G.apply(weights_init)
D.apply(weights_init)

opt_G = torch.optim.Adam(G.parameters(), lr=LR, betas=BETAS)
opt_D = torch.optim.Adam(D.parameters(), lr=LR, betas=BETAS)
criterion = nn.BCELoss()

fixed_noise = torch.randn(64, LATENT_DIM, 1, 1, device=device)

# --- Training loop ---
for epoch in range(NUM_EPOCHS):
    for i, (real_imgs, _) in enumerate(dataloader):
        real_imgs = real_imgs.to(device)
        batch_size = real_imgs.size(0)
        real_labels = torch.ones(batch_size, 1, device=device)
        fake_labels = torch.zeros(batch_size, 1, device=device)

        # (1) Update discriminator
        noise = torch.randn(batch_size, LATENT_DIM, 1, 1, device=device)
        fake_imgs = G(noise).detach()

        loss_D_real = criterion(D(real_imgs), real_labels)
        loss_D_fake = criterion(D(fake_imgs), fake_labels)
        loss_D = loss_D_real + loss_D_fake

        opt_D.zero_grad()
        loss_D.backward()
        opt_D.step()

        # (2) Update generator
        noise = torch.randn(batch_size, LATENT_DIM, 1, 1, device=device)
        fake_imgs = G(noise)
        loss_G = criterion(D(fake_imgs), real_labels)  # non-saturating

        opt_G.zero_grad()
        loss_G.backward()
        opt_G.step()

    print(f"Epoch [{epoch+1}/{NUM_EPOCHS}]  "
          f"Loss_D: {loss_D.item():.4f}  Loss_G: {loss_G.item():.4f}")
\end{minted}
\end{codeblock}

\begin{codeoutput}
\begin{verbatim}
Epoch [1/25]   Loss_D: 0.5823  Loss_G: 1.8932
Epoch [2/25]   Loss_D: 0.4517  Loss_G: 2.1045
...
Epoch [10/25]  Loss_D: 0.6102  Loss_G: 1.3287
...
Epoch [25/25]  Loss_D: 0.6731  Loss_G: 0.9845
\end{verbatim}
\end{codeoutput}

\section{Evaluation metrics}
\label{sec:gan-metrics}

\subsection{Inception Score (IS)}

\begin{keyformulas}[title=\textbf{Inception Score}]\index{Inception Score}\index{IS}
\begin{equation}
\text{IS} = \exp\!\Big(\E_{x \sim p_g}\!\big[\KL\!\big(p(y \mid x) \,\|\, p(y)\big)\big]\Big)
\end{equation}
where $p(y \mid x)$ is given by a pre-trained Inception network and $p(y) = \E_x[p(y \mid x)]$.

A high IS means: (1) the images are sharp ($p(y|x)$ concentrated) and (2) the images are diverse ($p(y)$ uniform).
\end{keyformulas}

\subsection{Fr\'echet Inception Distance (FID)}

\begin{keyformulas}[title=\textbf{Fr\'echet Inception Distance}]\index{FID}\index{Fr\'echet Inception Distance}
Features are extracted from the pool3 layer of the Inception network for real data ($\mu_r, \Sigma_r$) and generated data ($\mu_g, \Sigma_g$):
\begin{equation}
\text{FID} = \|\mu_r - \mu_g\|_2^2 + \text{Tr}\!\Big(\Sigma_r + \Sigma_g - 2\big(\Sigma_r \Sigma_g\big)^{1/2}\Big)
\label{eq:fid}
\end{equation}
A \textbf{low} FID indicates better quality and diversity.
\end{keyformulas}

\begin{remark}
FID is preferred over IS because it compares the generated distribution to the \emph{real} data, whereas IS only evaluates the internal properties of the generated images.
\end{remark}

\section{Ablation study}
\label{sec:gan-ablation}

\begin{table}[htbp]
\centering
\caption{DCGAN ablation study on Fashion-MNIST. FID measured after 25 epochs (lower = better).}
\label{tab:gan-ablation}
\begin{tabular}{lccc}
\toprule
\textbf{Configuration} & \textbf{Latent dim.} & \textbf{LR ($G$ / $D$)} & \textbf{FID} $\downarrow$ \\
\midrule
Base & 100 & $2\!\times\!10^{-4}$ / $2\!\times\!10^{-4}$ & 28.3 \\
\midrule
$z \in \R^{32}$ & 32 & $2\!\times\!10^{-4}$ / $2\!\times\!10^{-4}$ & 35.1 \\
$z \in \R^{64}$ & 64 & $2\!\times\!10^{-4}$ / $2\!\times\!10^{-4}$ & 30.2 \\
$z \in \R^{256}$ & 256 & $2\!\times\!10^{-4}$ / $2\!\times\!10^{-4}$ & 27.8 \\
\midrule
Higher $G$ LR & 100 & $5\!\times\!10^{-4}$ / $2\!\times\!10^{-4}$ & 32.5 \\
Higher $D$ LR & 100 & $2\!\times\!10^{-4}$ / $5\!\times\!10^{-4}$ & 41.7 \\
TTUR ($G < D$) & 100 & $1\!\times\!10^{-4}$ / $4\!\times\!10^{-4}$ & 25.9 \\
\midrule
$n_{\text{critic}}=1$ (standard) & 100 & $2\!\times\!10^{-4}$ / $2\!\times\!10^{-4}$ & 28.3 \\
$n_{\text{critic}}=3$ & 100 & $2\!\times\!10^{-4}$ / $2\!\times\!10^{-4}$ & 26.7 \\
$n_{\text{critic}}=5$ (WGAN-GP) & 100 & $1\!\times\!10^{-4}$ / $1\!\times\!10^{-4}$ & 23.4 \\
\bottomrule
\end{tabular}
\end{table}

\begin{remark}
Key findings: (1) the latent dimension has a moderate impact, but too small it limits diversity; (2) a higher learning rate for $D$ (TTUR --- Two Time-scale Update Rule)\index{TTUR} improves FID; (3) WGAN-GP with $n_{\text{critic}}=5$ achieves the best performance.
\end{remark}

\section{State of the art and connections}
\label{sec:gan-sota}

\begin{sota}[title=\textbf{Beyond classical GANs}]
\begin{itemize}
    \item \textbf{StyleGAN}\index{StyleGAN} (Karras et al., 2019--2021): introduces a mapping network $z \to w$ and style injection via \emph{Adaptive Instance Normalization} (AdaIN). StyleGAN3 eliminates texture artifacts through strict translation equivariance.

    \item \textbf{Diffusion models}\index{diffusion models} (Ho et al., 2020; Dhariwal and Nichol, 2021): denoising diffusion models have largely surpassed GANs in terms of FID on ImageNet, while offering superior diversity and more stable training. See DALL-E 2, Stable Diffusion, Imagen.

    \item \textbf{Consistency Models}\index{Consistency Models} (Song et al., 2023): distill diffusion models to enable single-step generation, combining diffusion quality with GAN speed.

    \item GANs remain preferred in applications requiring ultra-fast generation (real-time super-resolution, video synthesis).
\end{itemize}
\end{sota}

\section{Exercises}
\label{sec:gan-exercises}

\begin{exercise}[Verifying the optimal discriminator]
\label{exo:optimal-disc}
\textbf{Difficulty: 1/3}
Let $p_{\text{data}} = \mathcal{N}(3, 1)$ and $p_g = \mathcal{N}(0, 1)$. Analytically compute $D^*(x)$ using Equation~\eqref{eq:optimal-D}. Plot $D^*(x)$ for $x \in [-5, 8]$. Comment on the shape of the curve.
\end{exercise}

\begin{exercise}[WGAN: comparison with standard GAN]
\label{exo:wgan-compare}
\textbf{Difficulty: 2/3}
Consider two distributions: $p_{\text{data}} = \delta_0$ (Dirac mass at 0) and $p_g = \delta_\theta$.
\begin{enumerate}
    \item Compute $\text{JSD}(p_{\text{data}} \| p_g)$ for $\theta \neq 0$ and $\theta = 0$.
    \item Compute $W(p_{\text{data}}, p_g) = |\theta|$.
    \item Deduce why $W$ provides a useful gradient for $G$ while JSD does not.
\end{enumerate}
\end{exercise}

\begin{exercise}[WGAN-GP implementation]
\label{exo:wgan-gp-impl}
\textbf{Difficulty: 2/3}
Modify the DCGAN code from Section~\ref{sec:dcgan-pytorch} to:
\begin{enumerate}
    \item Remove the sigmoid from the discriminator (``critic'').
    \item Implement the gradient penalty (Equation~\eqref{eq:wgan-gp}).
    \item Train the critic $n_{\text{critic}}=5$ times per generator iteration.
    \item Compare the loss curves and visual quality with the standard GAN.
\end{enumerate}
\end{exercise}

\begin{exercise}[Conditional GAN for targeted generation]
\label{exo:cgan-impl}
\textbf{Difficulty: 3/3}
Implement a cGAN\index{cGAN} on Fashion-MNIST that takes a class label as input and generates an image of the corresponding category.

\emph{Hint}: concatenate a \emph{one-hot encoding} of the label to the latent vector for $G$, and add a label \emph{embedding} as an additional channel for $D$.
\end{exercise}

\begin{exercise}[Theoretical analysis --- GAN convergence]
\label{exo:gan-convergence}
\textbf{Difficulty: 3/3}
\begin{enumerate}
    \item Show that if $D$ and $G$ are updated simultaneously by gradient descent on $V(D, G)$, the dynamical system is \emph{not} a gradient game (the gradients are not conservative). What does this imply for convergence?
    \item Show that for the standard GAN in dimension 1, with $p_{\text{data}} = \mathcal{N}(0, 1)$ and $G(z) = \mu + \sigma z$ ($z \sim \mathcal{N}(0,1)$), the fixed point $\mu=0, \sigma=1$ is not a stable attractor of the simultaneous gradient dynamical system.
\end{enumerate}
\end{exercise}

\begin{exercise}[FID computation {\normalfont (project)}]
\label{exo:fid-compute}
\textbf{Difficulty: 3/3}
\begin{enumerate}
    \item Implement FID computation using a pre-trained network (e.g., InceptionV3 or a simpler network for $28 \times 28$ images).
    \item Compute the FID of your DCGAN at different training epochs.
    \item Plot the evolution of FID as a function of epoch and comment.
\end{enumerate}
\end{exercise}
